\section{Methodology: HyperMerge}
\label{sec:methodology}

We present \textbf{HyperMerge} (\textbf{Hyperbolic Space Activation Routing and Fusion}), a dynamic model ensembling system built entirely upon non-Euclidean geometry. Below, we detail our geometric substrate, activation mappings, algebraic primitives, and our two core algorithms: Hyperbolic Centroid Alignment (HCA) and Beltrami-Klein Symmetric Blending (BKSB).

\subsection{Poincaré Ball Model of Hyperbolic Space}
We select the Poincaré Ball model as our geometric workspace. Let $\mathbb{D}_c^D = \{\mathbf{x} \in \mathbb{R}^D : \|\mathbf{x}\|_2 < 1 / \sqrt{c}\}$ be the open $D$-dimensional ball with constant negative curvature $-c$ (where $c > 0$). The Poincaré Ball is a Riemannian manifold equipped with the metric tensor $g_{\mathbf{x}} = \lambda_{\mathbf{x}}^2 g_{\text{Eucl}}$, where $g_{\text{Eucl}} = \mathbf{I}_D$ is the standard Euclidean metric, and
\begin{equation}
\lambda_{\mathbf{x}} = \frac{2}{1 - c\|\mathbf{x}\|_2^2}
\end{equation}
is the conformal factor. As a point $\mathbf{x}$ approaches the boundary of the ball ($\|\mathbf{x}\|_2 \to 1/\sqrt{c}$), the conformal factor $\lambda_{\mathbf{x}}$ diverges to infinity. This metric structure results in exponential volume growth, providing a natural substrate for embedding deep, representation-level hierarchies and segregating disjoint task manifolds.

\subsection{Differentiable Mapping Functions}
To bridge our Euclidean neural network layers with the hyperbolic workspace, we utilize the exponential and logarithmic mapping functions defined at the origin $\mathbf{0}$. For any Euclidean activation vector $h \in \mathbb{R}^D \setminus \{\mathbf{0}\}$, the exponential map $\exp_{\mathbf{0}}^c: \mathbb{R}^D \to \mathbb{D}_c^D$ projects it into the Poincaré Ball:
\begin{equation}
\mathbf{z} = \exp_{\mathbf{0}}^c(h) = \tanh(\sqrt{c} \|h\|_2) \frac{h}{\sqrt{c} \|h\|_2},
\end{equation}
and $\exp_{\mathbf{0}}^c(\mathbf{0}) = \mathbf{0}$. 
Conversely, to project a hyperbolic state $\mathbf{z} \in \mathbb{D}_c^D \setminus \{\mathbf{0}\}$ back into flat Euclidean space, we define the logarithmic map $\log_{\mathbf{0}}^c: \mathbb{D}_c^D \to \mathbb{R}^D$:
\begin{equation}
h = \log_{\mathbf{0}}^c(\mathbf{z}) = \text{artanh}(\sqrt{c} \|\mathbf{z}\|_2) \frac{\mathbf{z}}{\sqrt{c} \|\mathbf{z}\|_2},
\end{equation}
and $\log_{\mathbf{0}}^c(\mathbf{0}) = \mathbf{0}$. Both mapping functions are closed-form and fully differentiable, allowing seamless integration into standard backpropagation pipelines.

\subsection{Möbius Algebraic Primitives}
To perform ensembling and scaling directly in hyperbolic space, we leverage the formalism of Möbius gyrovector spaces \cite{ganea2018hyperbolic}. 
For any two points $\mathbf{x}, \mathbf{y} \in \mathbb{D}_c^D$, the non-linear \emph{Möbius addition} ($\oplus_c$) is formulated as:
\begin{equation}
\mathbf{x} \oplus_c \mathbf{y} = \frac{(1 + 2c\langle\mathbf{x}, \mathbf{y}\rangle + c\|\mathbf{y}\|_2^2)\mathbf{x} + (1 - c\|\mathbf{x}\|_2^2)\mathbf{y}}{1 + 2c\langle\mathbf{x}, \mathbf{y}\rangle + c^2\|\mathbf{x}\|_2^2\|\mathbf{y}\|_2^2}.
\end{equation}
For any scalar scaling factor $r \in \mathbb{R}$ and point $\mathbf{x} \in \mathbb{D}_c^D \setminus \{\mathbf{0}\}$, the \emph{Möbius scalar multiplication} ($\otimes_c$) is defined as:
\begin{equation}
r \otimes_c \mathbf{x} = \frac{1}{\sqrt{c}} \tanh\left( r \cdot \text{artanh}(\sqrt{c} \|\mathbf{x}\|_2) \right) \frac{\mathbf{x}}{\|\mathbf{x}\|_2},
\end{equation}
and $r \otimes_c \mathbf{0} = \mathbf{0}$.

\subsection{Hyperbolic Distance and Routing}
The geodesic distance between any two points $\mathbf{x}, \mathbf{y} \in \mathbb{D}_c^D$ is given by:
\begin{equation}
d_{\mathbb{D}}^c(\mathbf{x}, \mathbf{y}) = \frac{2}{\sqrt{c}} \text{artanh}\left( \sqrt{c} \|\!-\mathbf{x} \oplus_c \mathbf{y}\|_2 \right).
\end{equation}
For an online sample $b$ with Layer 0 embedding representation $z_b^{\text{embed}}$, we project it to the Poincaré Ball: $\mathbf{z}_b = \exp_{\mathbf{0}}^c(z_b^{\text{embed}})$. The hyperbolic distance to each pre-computed task centroid reference $\boldsymbol{\mu}_k$ is $D_{k, b} = d_{\mathbb{D}}^c(\mathbf{z}_b, \boldsymbol{\mu}_k)$. To compute dynamic, sample-wise ensembling coefficients $\alpha_{k, b}$, we perform temperature-scaled Softmax over the negative distance coordinates:
\begin{equation}
\alpha_{k, b} = \frac{\exp(-D_{k, b} / \tau)}{\sum_{j=1}^K \exp(-D_{j, b} / \tau)},
\label{eq:routing}
\end{equation}
where $\tau > 0$ is the routing temperature.

\subsection{Hyperbolic Centroid Alignment (HCA)}
Computing the centroid (the Fréchet mean) of a calibration set $\mathcal{C}_k = \{z_{1}, \dots, z_{N}\}$ for task $k$ in the Poincaré metric is highly non-trivial. However, we can obtain a closed-form solution by mapping the Poincaré coordinates to the Beltrami-Klein model $\mathbb{K}_c^D$ of hyperbolic space, where geodesics are straight lines.
First, we project each calibration embedding to Poincaré Ball coordinates: $\mathbf{z}_s = \exp_{\mathbf{0}}^c(z_s)$. Then, we map Poincaré coordinates to Klein coordinates:
\begin{equation}
\mathbf{z}_s^{Klein} = \frac{2 \mathbf{z}_s}{1 + c \|\mathbf{z}_s\|_2^2}.
\end{equation}
In Klein space, the barycenter corresponds exactly to the \emph{Einstein midpoint}, which is a weighted average of the coordinates weighted by their Lorentz factors $\gamma_s$:
\begin{equation}
\boldsymbol{\mu}_k^{Klein} = \frac{\sum_{s \in \mathcal{C}_k} \gamma_s \mathbf{z}_s^{Klein}}{\sum_{s \in \mathcal{C}_k} \gamma_s}, \quad \text{where } \gamma_s = \frac{1}{\sqrt{1 - c \|\mathbf{z}_s^{Klein}\|_2^2}}.
\end{equation}
Finally, we project the Einstein midpoint back into the Poincaré Ball coordinates to obtain the mathematically optimal task-reference centroid $\boldsymbol{\mu}_k$:
\begin{equation}
\boldsymbol{\mu}_k = \frac{\boldsymbol{\mu}_k^{Klein}}{1 + \sqrt{1 - c \|\boldsymbol{\mu}_k^{Klein}\|_2^2}}.
\end{equation}

\subsection{Beltrami-Klein Symmetric Blending (BKSB)}
At layer $l$, let $h_b^{(l-1)} \in \mathbb{R}^D$ be the activation propagated from the previous layer. The base model's output is $h_{base, b}^{(l)} = h_b^{(l-1)} W_{\text{base}}^{(l)} \in \mathbb{R}^D$, and each task expert's low-rank adapter update is $E_{k, b}^{(l)} = h_b^{(l-1)} A_k^{(l)} B_k^{(l)} \in \mathbb{R}^D$.
To perform ensembling without flat-space crowding, we project each expert update into the Poincaré Ball:
\begin{equation}
\mathbf{v}_{k, b} = \exp_{\mathbf{0}}^c\left( E_{k, b}^{(l)} \right).
\end{equation}
Because Möbius addition ($\oplus_c$) is non-associative and non-commutative, a sequential accumulation would introduce an arbitrary task-indexing order bias. Furthermore, scaling expert updates prior to projection via Möbius scalar multiplication and then performing a weighted sum in Klein space (as described in some ad-hoc ensembling schemes) leads to a severe \emph{double-weighting flaw} where the updates are scaled by $\alpha_{k, b}^2$, artificially shrinking their norms and degrading ensembling quality.

To achieve a mathematically consistent and completely permutation-invariant ensembling with a single application of routing weights, we project the unscaled Poincaré updates directly to the Beltrami-Klein model $\mathbb{K}_c^D$ where geodesics are straight lines:
\begin{equation}
\mathbf{w}_{k, b} = \frac{2 \mathbf{v}_{k, b}}{1 + c \|\mathbf{v}_{k, b}\|_2^2}.
\end{equation}
Once in the Beltrami-Klein space, we compute the mathematically optimal hyperbolic barycenter (corresponding to the true Fréchet mean on the manifold). This barycenter corresponds exactly to the Lorentz-weighted \emph{Einstein midpoint}, where each Klein coordinate $\mathbf{w}_{k, b}$ is scaled by its dynamic routing weight $\alpha_{k, b}$ and its Lorentz factor $\gamma_{k, b}$:
\begin{equation}
\mathbf{w}_{\text{merged}, b} = \frac{\sum_{k=1}^K \alpha_{k, b} \gamma_{k, b} \mathbf{w}_{k, b}}{\sum_{k=1}^K \alpha_{k, b} \gamma_{k, b}}, \quad \text{where } \gamma_{k, b} = \frac{1}{\sqrt{1 - c \|\mathbf{w}_{k, b}\|_2^2}}.
\end{equation}
Since the Lorentz-weighted Einstein midpoint is associative and commutative, this formulation is completely permutation-invariant and applies the routing weights exactly once, preserving the correct magnitude of expert updates.

We then map the merged Klein state back to Poincaré Ball coordinates:
\begin{equation}
\mathbf{v}_{\text{merged}, b} = \frac{\mathbf{w}_{\text{merged}, b}}{1 + \sqrt{1 - c \|\mathbf{w}_{\text{merged}, b}\|_2^2}}.
\end{equation}
Finally, we map the merged update back to flat Euclidean space via the logarithmic map:
\begin{equation}
E_{\text{merged}, b}^{(l)} = \log_{\mathbf{0}}^c\left( \mathbf{v}_{\text{merged}, b} \right).
\end{equation}
The final activation state propagated to the next layer is:
\begin{equation}
h_b^{(l)} = h_{base, b}^{(l)} + E_{\text{merged}, b}^{(l)}.
\end{equation}
By keeping the base model in Euclidean space and performing non-linear ensembling of adapter updates inside the Poincaré Ball via Klein symmetric weighted averages, HyperMerge achieves a highly expressive, non-linear activation blend that is completely symmetric with respect to task indices and segregates conflicting task trajectories towards the hyperbolic boundary.

\subsection{Hyperbolic Out-of-Distribution Rejection (HOR)}
To secure edge deployments against malicious or out-of-distribution (OOD) tasks, HyperMerge implements a non-parametric outlier detector. For each incoming sample $b$, we evaluate its minimum hyperbolic distance to any known task centroid:
\begin{equation}
d_{\text{min}, b} = \min_{k \in \{1..K\}} d_{\mathbb{D}}^c(\mathbf{z}_b, \boldsymbol{\mu}_k).
\end{equation}
If $d_{\text{min}, b} > \gamma_{\text{OOD}}$, where $\gamma_{\text{OOD}}$ is a distance threshold, the query is rejected as out-of-distribution, protecting downstream experts from processing toxic or irrelevant queries.

\subsection{Analysis of Projection Distortion and Geometric Consistency}
HyperMerge operates as a hybrid Euclidean-hyperbolic system, where the base model remains in flat Euclidean space $\mathbb{R}^D$ while activation-space expert updates are projected to the Poincaré Ball $\mathbb{D}_c^D$ for non-linear ensembling. We justify the geometric consistency and analyze the distortion of this hybrid design below.

\textbf{Tangent-Space Approximation and Consistency:} We treat the Euclidean representation space $\mathbb{R}^D$ of the pre-trained base model as the tangent space $T_{\mathbf{0}}\mathbb{D}_c^D$ at the origin of the hyperbolic manifold. The zero vector represents the undisturbed base model state, and the low-rank adapter updates $E_{k, b}^{(l)}$ represent directional displacement vectors residing in this tangent space. Projecting these displacements via the exponential map $\exp_{\mathbf{0}}^c$ is equivalent to mapping tangent vectors onto the manifold. The Beltrami-Klein Symmetric Blending (BKSB) then aggregates these points on the manifold, utilizing its negative curvature to prevent representation crowding. Finally, mapping the merged state back to the origin's tangent space via the logarithmic map $\log_{\mathbf{0}}^c$ yields a consistent Euclidean displacement vector $E_{\text{merged}, b}^{(l)}$ which can be linearly added to the Euclidean base activation $h_{base, b}^{(l)}$. This formulation is mathematically analogous to the retraction operations used in Riemannian optimization, guaranteeing geometric consistency.

\textbf{Information Loss and Distortion Bounds:} Because $\exp_{\mathbf{0}}^c$ and $\log_{\mathbf{0}}^c$ are non-linear, repeatedly mapping activations back and forth introduces radial compression and expansion distortion. Let $h \in T_{\mathbf{0}}\mathbb{D}_c^D$ be an update vector, and let its projection be $\mathbf{z} = \exp_{\mathbf{0}}^c(h)$. By performing a Taylor expansion of the mapping functions, we can bound the projection distortion. Since $\tanh(x) = x - \frac{1}{3}x^3 + O(x^5)$, we have:
\begin{equation}
\mathbf{z} = \left( \sqrt{c}\|h\|_2 - \frac{c^{3/2}}{3}\|h\|_2^3 + O(c^{5/2}\|h\|_2^5) \right) \frac{h}{\sqrt{c}\|h\|_2} = h - \frac{c}{3}\|h\|_2^2 h + O(c^2 \|h\|_2^4 h).
\end{equation}
The projection distortion $\delta = \|\mathbf{z} - h\|_2$ is bounded by:
\begin{equation}
\delta \leq \frac{c}{3}\|h\|_2^3 + O(c^2 \|h\|_2^5).
\end{equation}
For Low-Rank Adapters (LoRA), the updates $E_{k, b}^{(l)}$ are typically highly low-rank and have small norms ($\|E\|_2 \ll 1$). Consequently, under a moderate curvature $c = 0.1$, the cubic dependence on $\|h\|_2$ ensures that the projection distortion is negligible ($\delta \approx 10^{-4}$), meaning the hybrid mapping preserves almost all representation information while reaping the massive benefits of hyperbolic separation.

\subsection{Numerical Stability and Safeguards}
In hyperbolic neural networks, division by conformal factors and coordinate transformations can be highly sensitive to numerical precision, particularly near the boundary of the Poincaré Ball ($\|\mathbf{x}\|_2 \to 1/\sqrt{c}$) where the conformal factor $\lambda_{\mathbf{x}} \to \infty$. To prevent numerical overflow, imaginary square roots, or NaN errors, HyperMerge implements strict numerical safeguards in both Poincaré and Klein coordinates.

For any hyperbolic point $\mathbf{x} \in \mathbb{D}_c^D$ or its Klein counterpart $\mathbf{w} \in \mathbb{K}_c^D$, we enforce a strict bounding radius by clamping their Euclidean norms to a maximum value of $\frac{1}{\sqrt{c}} - \epsilon$, where $\epsilon = 10^{-7}$ represents a small precision guard:
\begin{equation}
\text{clip}(\mathbf{x}) = \min\left(1, \frac{1/\sqrt{c} - \epsilon}{\|\mathbf{x}\|_2}\right) \mathbf{x}.
\end{equation}
This clamping is applied: (1) immediately after the exponential map projection to ensure all projected activations are strictly inside the Poincaré Ball; (2) during coordinate mapping from Poincaré to Klein space; and (3) when computing the Lorentz factor $\gamma = \frac{1}{\sqrt{1 - c \|\mathbf{w}\|_2^2}}$, where the term under the square root is guaranteed to be strictly greater than $2\sqrt{c}\epsilon - c\epsilon^2 > 0$. This ensures the complete elimination of NaNs and guarantees numerical stability during deep, multi-layer propagation.
